\subsection{Relation to measures}\label{sec:measure}

Up to here this paper has treated integration without using measures.
In this subsection we state the correspondence between integration systems and
measure theory.
The main results of this paper do not use this correspondence, but we record that
the notion of an integration system does recover measure theory.
Everything below is known, and we refer to the literature for proofs.

\subsubsection{From measures to integration systems}

This direction is the easy one.

\begin{proposition}\label{prop:measure-gives-integral-system}
  Let $(X,\Sigma,\mu)$ be a measure space, let $L$ be the set of all
  $\Sigma$-measurable simple functions that are integrable with respect to
  $\mu$, and let $I(f):=\int_Xf\,d\mu$.
  Then $(X,L,I)$ is an integration system.
\end{proposition}

\noindent
Positivity is immediate, and verifying the continuity axiom uses
$\sigma$-additivity of $\mu$.
For a proof see \cite{BM}, \S7.
The set of step functions in Example~\ref{ex:step-functions} is a subsystem of
the set of all simple functions with respect to Lebesgue measure.
Finite unions of half-open intervals do not form a $\sigma$-algebra, so it is not
itself a special case of this proposition.

\subsubsection{The Stone condition}

The converse direction, that of extracting a measure from an integration system,
needs one condition.

\begin{definition}\label{def:stone-condition}
  We say that an integration system $(X,L,I)$ on a set $X$ satisfies
  the \textbf{Stone condition} if
  \[
    f\wedge1\in L
  \]
  holds for every $f\in L$. Here $1$ is the constant function $1$ on $X$.
\end{definition}

\noindent
Note that $L$ itself need not contain $1$.
What is required is that the function obtained by truncating $f$ at $1$ belong to
$L$ again.
This condition is the standard one, and a Riesz space satisfying it is called a
\textbf{Stone vector lattice}.

This condition is needed because we want to define the measure of a set $A$ with
$1_A\in L^1$ by
\[
  \mu(A):=I^1(1_A) .
\]
If $L^1$ contains no indicator function at all, this definition cannot even be
started.

\begin{remark}\label{rem:mu-only-on-finite}
  The formula above is meaningful only for sets with $1_A\in L^1$.
  It cannot be used for general measurable sets, and in particular not for sets
  of infinite measure.

  Looking at Example~\ref{ex:counting} ($X=\mathbb{N}$, the functions of finite
  support, and summation), $L^1$ is the set $\ell^1$ of all absolutely summable
  sequences and the measure one should obtain is the counting measure.
  But $1_{\mathbb{N}}\notin\ell^1$, so $I^1(1_{\mathbb{N}})$ is undefined, while
  $\mu(\mathbb{N})=\infty$.

  To construct a general measure one must therefore start from the sets with
  $1_A\in L^1$ and pass through an upper integral or an outer measure.
  Theorem~\ref{thm:daniell-stone} below states only the result of that
  construction.
\end{remark}

\begin{example}\label{ex:stone-holds}
  Both Example~\ref{ex:counting} and Example~\ref{ex:step-functions} satisfy the
  Stone condition.
  Indeed, truncating a function of finite support at $1$ again has finite
  support, and truncating a step function at $1$ is again a step function.
\end{example}

\begin{example}[an integration system not satisfying the Stone condition]\label{ex:stone-fails}
  Let $X=\{1,2\}$ and put
  \[
    L:=\{(a,2a)\mid a\in\mathbb{R}\}\subset\mathbb{R}^X,
    \qquad
    I(a,2a):=a .
  \]
  Then $L$ is the one-dimensional subspace spanned by $(1,2)$, and
  $|(a,2a)|=|a|\cdot(1,2)\in L$, so $L$ is a Riesz space on $X$.
  Moreover $I$ is clearly positive linear, and continuity holds because $X$ is
  finite.
  Hence $(X,L,I)$ is an integration system.

  But $f=(1,2)$ gives $f\wedge1=(1,1)\notin L$, so the Stone condition fails.

  In this integration system $L^1=L$.
  Indeed, the dominating functions occurring in a series representation have the
  form $\varphi_n=(c_n,2c_n)$ with $c_n\ge0$, and if
  $\sum_nI(\varphi_n)=\sum_nc_n<+\infty$ then $\sum_n\varphi_n(x)<+\infty$ holds
  at both $x=1$ and $x=2$.
  Hence condition (3) of a series representation is imposed on all of $X$, and
  $f$ is confined to being the sum of an absolutely convergent series of
  elements of $L$, that is, an element of $L$.

  Consequently $1_{\{1\}}=(1,0)\notin L^1$, and $\mu(\{1\}):=I^1(1_{\{1\}})$
  cannot be defined.
\end{example}

\subsubsection{The Daniell--Stone representation theorem}

\begin{theorem}[Daniell--Stone]\label{thm:daniell-stone}
  Let $(X,L,I)$ be an integration system satisfying the Stone condition.
  Then there exist a $\sigma$-algebra $\Sigma$ on $X$ and a measure $\mu$ such
  that every $f\in L$ is $\mu$-integrable and
  \[
    I(f)=\int_Xf\,d\mu .
  \]
\end{theorem}

\noindent
For a proof see \cite{Stone}, \cite{Royden} Chapter 16, or \cite{CFW}.
Chapter 16 of \cite{Royden} is entitled ``The Daniell Integral''; its \S2 is the
extension theorem, \S3 uniqueness, and \S4 measurability and measure.
Here $\Sigma$ is chosen so that the functions belonging to $L^1$ are measurable,
and $\mu$ is constructed through an outer measure, as in
Remark~\ref{rem:mu-only-on-finite}, so as to satisfy $\mu(A)=I^1(1_A)$ for the
sets $A$ with $1_A\in L^1$.

\begin{example}
  Applying this theorem to Example~\ref{ex:step-functions}, the measure obtained
  is Lebesgue measure on $[0,1)$ and the Daniell extension $L^1$ coincides with
  the set of all Lebesgue integrable functions.
  This is how the fact stated in Example~\ref{ex:lebesgue} looks from the side
  of measure theory.
\end{example}

\begin{remark}
  This paper does not impose the Stone condition.
  The results of Section~\ref{chap:functoriality} and
  Section~\ref{chap:operator} hold without it.
  Indeed the integration system of Example~\ref{ex:stone-fails} is an object of
  $\mathbf{Int}$ as well, and is moreover Daniell closed.

  One cannot, however, conclude from this that ``the framework of this paper is
  wider than measure theory''.
  The integral of Example~\ref{ex:stone-fails} can be written as
  $I(a,2a)=\int_Xf\,d\mu$ by taking $\mu(\{1\})=1$ and $\mu(\{2\})=0$
  (and such a $\mu$ is moreover not unique: any $\mu$ with
  $\mu(\{1\})+2\mu(\{2\})=1$ will do).
  What that example shows is that without the Stone condition one
  \textbf{cannot canonically recover a measure from inside the Daniell extension
  using indicator functions}, and that \textbf{the representing measure is not
  uniquely determined}.

  Put precisely, this paper does not assume the Stone condition and does not take
  representation by a measure as a premise of the theory.
\end{remark}

\begin{problem}\label{prob:stone-measure-spaces}
  Is every Daniell closed integration system satisfying the Stone condition
  realized as the $L^1$ of a measure space, and does this correspondence lift to
  an equivalence of categories?
  That is, is the full subcategory of $\mathbf{DInt}$ consisting of the objects
  that satisfy the Stone condition equivalent to a suitable category of measure
  spaces and measure-preserving maps?
\end{problem}

\noindent
This question was raised by T.~Fritz in personal communication (2026).
It does not hold, however, in the form of an equivalence with the category of
measurable spaces.
A morphism of integration systems preserves the integral, $J(F^*f)=I(f)$, so the
integral itself is data that distinguishes objects up to isomorphism.
Indeed, let $X=\{*\}$, $L=\mathbb{R}$ and $I_c(a):=ca$ for $c>0$.
Since $L^1\subset\mathbb{R}^X=\mathbb{R}=L$, each of these is a Daniell closed
integration system satisfying the Stone condition.
But $(X,\mathbb{R},I_c)\not\cong(X,\mathbb{R},I_{c'})$ for $c\ne c'$, while the
underlying measurable space is a one-point space in every case.
What should be asked about is therefore measure spaces, not measurable spaces.
Moreover, by Theorem~\ref{thm:complete-intrinsic} a Daniell closed integration
system is closed under modification on null sets, so the measures to be matched
are expected to be the complete ones.
